\documentclass[a4paper, 14pt]{extarticle}

% Поля
%--------------------------------------
\usepackage{geometry}
\geometry{a4paper,tmargin=2cm,bmargin=2cm,lmargin=3cm,rmargin=1cm}
%--------------------------------------

\usepackage{float}
%Russian-specific packages
%--------------------------------------
\usepackage[T2A]{fontenc}
\usepackage[utf8]{inputenc} 
\usepackage[english, main=russian]{babel}
%--------------------------------------

\usepackage{textcomp}

% Красная строка
%--------------------------------------
\usepackage{indentfirst}               
%--------------------------------------             
\usepackage[most]{tcolorbox}
\tcbuselibrary{listings,breakable}
\newtcblisting{breakablelisting}{
  listing only,
  breakable,
  enhanced,
  frame hidden,
  colback=white,
  listing options={
    mathescape=false,
    texcl=false,
    inputencoding=utf8,
    literate={\$}{{\$}}1
  }
}


%Graphics
%--------------------------------------
\usepackage{graphicx}
\usepackage{wrapfig}
%--------------------------------------
\graphicspath{{./lab5/}}
% Полуторный интервал
%--------------------------------------
\linespread{1.3}                    
%--------------------------------------

%Выравнивание и переносы
%--------------------------------------
% Избавляемся от переполнений
\sloppy
% Запрещаем разрыв страницы после первой строки абзаца
\clubpenalty=10000
% Запрещаем разрыв страницы после последней строки абзаца
\widowpenalty=10000
%--------------------------------------

%Списки
\usepackage{enumitem}

%Подписи
\usepackage{caption} 

%Гиперссылки
\usepackage{hyperref}

\hypersetup {
	unicode=true
}

%Рисунки
%--------------------------------------
\DeclareCaptionLabelSeparator*{emdash}{~--- }
\captionsetup[figure]{labelsep=emdash,font=onehalfspacing,position=bottom}
%--------------------------------------

\usepackage{tempora}

%Листинги
%--------------------------------------
\usepackage{listings}
\lstset{
  basicstyle=\ttfamily\footnotesize, 
  %basicstyle=\footnotesize\AnkaCoder,        % the size of the fonts that are used for the code
  breakatwhitespace=false,         % sets if automatic breaks shoulbd only happen at whitespace
  breaklines=true,                 % sets automatic line breaking
  captionpos=t,                    % sets the caption-position to bottom
  inputencoding=utf8,
  frame=single,                    % adds a frame around the code
  keepspaces=true,                 % keeps spaces in text, useful for keeping indentation of code (possibly needs columns=flexible)
  keywordstyle=\bf,       % keyword style
  numbers=left,                    % where to put the line-numbers; possible values are (none, left, right)
  numbersep=5pt,                   % how far the line-numbers are from the code
  xleftmargin=25pt,
  xrightmargin=25pt,
  showspaces=false,                % show spaces everywhere adding particular underscores; it overrides 'showstringspaces'
  showstringspaces=false,          % underline spaces within strings only
  showtabs=false,                  % show tabs within strings adding particular underscores
  stepnumber=1,                    % the step between two line-numbers. If it's 1, each line will be numbered
  tabsize=2,                       % sets default tabsize to 8 spaces
  title=\lstname                   % show the filename of files included with \lstinputlisting; also try caption instead of title
}
%--------------------------------------

%%% Математические пакеты %%%
%--------------------------------------
\usepackage{amsthm,amsfonts,amsmath,amssymb,amscd}  % Математические дополнения от AMS
\usepackage{mathtools}                              % Добавляет окружение multlined
\usepackage[perpage]{footmisc}
%--------------------------------------

%--------------------------------------
%			НАЧАЛО ДОКУМЕНТА
%--------------------------------------

\begin{document}

%--------------------------------------
%			ТИТУЛЬНЫЙ ЛИСТ
%--------------------------------------
\begin{titlepage}
\thispagestyle{empty}
\newpage


%Шапка титульного листа
%--------------------------------------
\vspace*{-60pt}
\hspace{-65pt}
\begin{minipage}{0.3\textwidth}
\hspace*{-20pt}\centering
\includegraphics[width=\textwidth]{emblem}
\end{minipage}
\begin{minipage}{0.67\textwidth}\small \textbf{
\vspace*{-0.7ex}
\hspace*{-6pt}\centerline{Министерство науки и высшего образования Российской Федерации}
\vspace*{-0.7ex}
\centerline{Федеральное государственное автономное образовательное учреждение }
\vspace*{-0.7ex}
\centerline{высшего образования}
\vspace*{-0.7ex}
\centerline{<<Московский государственный технический университет}
\vspace*{-0.7ex}
\centerline{имени Н.Э. Баумана}
\vspace*{-0.7ex}
\centerline{(национальный исследовательский университет)>>}
\vspace*{-0.7ex}
\centerline{(МГТУ им. Н.Э. Баумана)}}
\end{minipage}
%--------------------------------------

%Полосы
%--------------------------------------
\vspace{-25pt}
\hspace{-35pt}\rule{\textwidth}{2.3pt}

\vspace*{-20.3pt}
\hspace{-35pt}\rule{\textwidth}{0.4pt}
%--------------------------------------

\vspace{1.5ex}
\hspace{-35pt} \noindent \small ФАКУЛЬТЕТ\hspace{80pt} <<Информатика и системы управления>>

\vspace*{-16pt}
\hspace{47pt}\rule{0.83\textwidth}{0.4pt}

\vspace{0.5ex}
\hspace{-35pt} \noindent \small КАФЕДРА\hspace{50pt} <<Теоретическая информатика и компьютерные технологии>>

\vspace*{-16pt}
\hspace{30pt}\rule{0.866\textwidth}{0.4pt}
  
\vspace{11em}

\begin{center}
\Large {\bf Лабораторная работа № 6} \\ 
\large {\bf по курсу <<Методы оптимизации>>} \\
\large <<Освоение многопараметрических методов>> 
\end{center}\normalsize

\vspace{8em}


\begin{flushright}
  {Студент группы ИУ9-82Б Шемякин В.А. \hspace*{15pt}\\ 
  \vspace{2ex}
  Преподаватель Посевин Д. П.\hspace*{15pt}}
\end{flushright}

\bigskip

\vfill
 

\begin{center}
\textsl{Москва 2026}
\end{center}
\end{titlepage}
%--------------------------------------
%		КОНЕЦ ТИТУЛЬНОГО ЛИСТА
%--------------------------------------

\renewcommand{\ttdefault}{pcr}

\setlength{\tabcolsep}{3pt}
\newpage
\setcounter{page}{2}

\section{Задание}\label{Sect::task}
Необходимо реализовать простейший многопараметрический метод оптимизации для двух параметров.
\section{Результаты}\label{Sect::res}

Исходный код программы представлен в листинге 1.

\begin{breakablelisting}
using LinearAlgebra
using Plots
using Colors

rosenbrock(x) = sum(100 * (x[i+1] - x[i]^2)^2 + (x[i] - 1)^2 for i in 1:length(x)-1)
schwefel(x) = 418.9829 * length(x) - sum(xi -> xi * sin(sqrt(abs(xi))), x)
rastrigin(x) = 10 * length(x) + sum(xi -> xi^2 - 10 * cos(2π * xi), x)

function grad_rosenbrock(x)
    n = length(x)
    g = zeros(n)
    for i in 1:n
        if i == 1
            g[i] = -400 * x[1] * (x[2] - x[1]^2) + 2 * (x[1] - 1)
        elseif i == n
            g[i] = 200 * (x[n] - x[n-1]^2) + 2 * (x[n] - 1)
        else
            g[i] = 200 * (x[i] - x[i-1]^2) - 400 * x[i] * (x[i+1] - x[i]^2) + 2 * (x[i] - 1)
        end
    end
    return g
end

function grad_schwefel(x)
    g = similar(x)
    for i in eachindex(x)
        xi = x[i]
        s = sqrt(abs(xi)) + 1e-20
        g[i] = xi >= 0 ? -sin(s) - xi * cos(s) / (2 * s) : -sin(s) + xi * cos(s) / (2 * s)
    end
    return g
end

grad_rastrigin(x) = [2 * xi + 20π * sin(2π * xi) for xi in x]

function golden_section(f, a, b, tol=1e-8)
    phi = (sqrt(5) - 1) / 2
    c = b - phi * (b - a)
    d = a + phi * (b - a)
    while abs(c - d) > tol
        if f(c) < f(d)
            b = d
        else
            a = c
        end
        c = b - phi * (b - a)
        d = a + phi * (b - a)
    end
    return (a + b) / 2
end

function minimize_2d(f, l0_range, l1_range, max_inner=50)
    l0, l1 = (l0_range[1] + l0_range[2]) / 2, (l1_range[1] + l1_range[2]) / 2
    for _ in 1:max_inner
        l0 = golden_section(x -> f([x, l1]), l0_range[1], l0_range[2])
        l1 = golden_section(x -> f([l0, x]), l1_range[1], l1_range[2])
    end
    return [l0, l1]
end

function mille_contrell(f, grad_f, x0, eps; max_iter=10000, zero_lambda1=false)
    n = length(x0)
    x = copy(x0)
    x_prev = copy(x0)
    delta_x = zeros(n)
    traj = [copy(x)]
    iter_count = 0

    for k in 1:max_iter
        iter_count = k
        if k > 0 && k % (n + 1) == 0
            delta_x = zeros(n)
        end
        g = grad_f(x)
        if norm(g) < eps
            break
        end

        if zero_lambda1
            delta_x = zeros(n)
        end
        obj(lambda) = f(x - lambda[1] * g + lambda[2] * delta_x)
        r = max(0.1, min(1.0, norm(delta_x) / (norm(g) + 1e-10)))
        l0_range = (0.0, r * 2)
        l1_range = (-r, r)
        l_opt = zero_lambda1 ? [golden_section(t -> obj([t, 0.0]), l0_range[1], l0_range[2]), 0.0] : minimize_2d(obj, l0_range, l1_range)

        x_new = x - l_opt[1] * g + l_opt[2] * delta_x
        delta_x = x_new - x
        x_prev = x
        x = x_new
        push!(traj, copy(x))

        if norm(delta_x) < eps
            break
        end
    end

    return x, f(x), traj, iter_count
end

function draw_trajectory_2d(f, traj, title_str, filename)
    xs = [p[1] for p in traj]
    ys = [p[2] for p in traj]
    fv = [f(p) for p in traj]
    rng = max(maximum(fv) - minimum(fv), 1e-10)
    span_x = maximum(xs) - minimum(xs)
    span_y = maximum(ys) - minimum(ys)
    pad = 0.2 * max(span_x, span_y, 0.5)
    xlim = (minimum(xs) - pad, maximum(xs) + pad)
    ylim = (minimum(ys) - pad, maximum(ys) + pad)
    xg = range(xlim[1], xlim[2], length=80)
    yg = range(ylim[1], ylim[2], length=80)
    plt = contour(xg, yg, (x, y) -> f([x, y]), levels=25, color=cgrad(:viridis), colorbar=true, legend=false, xlabel="x_1", ylabel="x_2", title=title_str, xlims=xlim, ylims=ylim)
    for i in 1:length(traj)-1
        t = (fv[i] - minimum(fv)) / rng
        col = RGB(0.2 + 0.6 * (1 - t), 0.3 + 0.5 * (1 - t), 0.8 - 0.3 * t)
        plot!(plt, xs[i:i+1], ys[i:i+1], color=col, lw=3, label="")
    end
    scatter!(plt, [xs[1]], [ys[1]], color=:green, ms=8, marker=:circle, label="")
    scatter!(plt, [xs[end]], [ys[end]], color=:red, ms=8, marker=:circle, label="")
    savefig(plt, filename)
    display(plt)
end

eps = 1e-6

funcs = [
    ("Розенброк", rosenbrock, grad_rosenbrock, [0.0, 0.0]),
    ("Швефель", schwefel, grad_schwefel, [350.0, 400.0]),
    ("Растригин", rastrigin, grad_rastrigin, [3.0, 3.0])
]

results = []
println("Результаты")
for (name, f, grad_f, x0) in funcs
    x, fv, traj, it = mille_contrell(f, grad_f, x0, eps)
    println(name, ":")
    println("  x* = ", round.(x, digits=6))
    println("  f(x*) = ", round(fv, digits=10))
    println("  Итерации = ", it)
    push!(results, (name=name, f=f, traj=traj))
end

println("\n lambda_1 = 0")
for (name, f, grad_f, x0) in funcs
    x, fv, traj, it = mille_contrell(f, grad_f, x0, eps; zero_lambda1=true)
    println(name, ":")
    println("  x* = ", round.(x, digits=6))
    println("  f(x*) = ", round(fv, digits=10))
    println("  Итерации = ", it)
end

draw_trajectory_2d(results[1].f, results[1].traj, "Розенброк", "lab6_rosenbrock.png")
draw_trajectory_2d(results[2].f, results[2].traj, "Швефель", "lab6_schwefel.png")
draw_trajectory_2d(results[3].f, results[3].traj, "Растригин", "lab6_rastrigin.png")
\end{breakablelisting}

\newpage

Результат представлен на рисунках~\ref{fig:img1}--~\ref{fig:rastrigin}.

\begin{figure}[H]
  \centering
  \includegraphics[width=0.6\textwidth]{lab6/screen1.png}
  \caption{Результат}
  \label{fig:img1}
\end{figure}

\begin{figure}[H]
  \centering
  \includegraphics[width=0.6\textwidth]{lab6/screen2.png}
  \caption{Результат}
  \label{fig:schwefel}
\end{figure}

\begin{figure}[H]
  \centering
  \includegraphics[width=0.6\textwidth]{lab6/screen3.png}
  \caption{Результат}
  \label{fig:rastrigin}
\end{figure}

\end{document}
